\documentclass[11pt]{article}
\usepackage[margin=1in]{geometry}
\usepackage{amsmath, amssymb}
\usepackage{graphicx}
\usepackage{xcolor}
\usepackage{hyperref}

\begin{document}

\textbf{DSC 210, Fall 2026 --- Homework 3} \hfill Nicholas Wendt (200478549)

%%%%%%%%%%%%%%%%%%%%%%%%%%%%%%%%%%%%%%%%%%%%%%%%%%%%%%%%%%%%%%%%%%%%%%%%
\section*{Question 1: Singular value decomposition (25 points)}

Determine the SVD of the following matrices by hand calculation. Make sure to explicitly list $\Sigma$, $U$, and $V$ in your answer.
\begin{enumerate}
    \item[(a)] $\begin{bmatrix}
    6 & 0 \\ 0 & 2
    \end{bmatrix}$ (4 points)
    \item[(b)] $\begin{bmatrix}
    2 & 0 \\ 0 & 6
    \end{bmatrix}$ (4 points)
    \item[(c)] $\begin{bmatrix}
    0 & 4 \\ 0 & 0 \\ 0 & 0
    \end{bmatrix}$ (4 points)
    \item[(d)] $\begin{bmatrix}
    1 & 0 \\ 2 & 0
    \end{bmatrix}$ (4 points)
    \item[(e)] $\begin{bmatrix}
    2 & 2 \\ 3 & 1
    \end{bmatrix}$ (4 points)
    \item[(f)] Write Python code to compute SVD of the matrices in parts (a) to (e) and verify that the above results match what you got from code. (5 points)\\
    \textbf{Hint}: You can use the numpy library function.
\end{enumerate}

\textbf{Solution:}

\bigskip

%%%%%%%%%%%%%%%%%%%%%%%%%%%%%%%%%%%%%%%%%%%%%%%%%%%%%%%%%%%%%%%%%%%%%%%%
\newpage
\section*{Question 2: Eigenvalues and Eigenvectors (25 points)}

$$
\mathbf{A} = \begin{bmatrix}
    4 & 2\\
    1 & 3
\end{bmatrix}
$$

\begin{enumerate}
    \item[(a)] Write down the characteristic equation for matrix $\mathbf{A}$. Use the characteristic equation to solve the eigenvalues and normalized eigenvectors of matrix $\mathbf{A}$. (10 points)
    \item[(b)] Write python code to verify your answers for the eigenvalues and normalized eigenvectors from Part (a). \textbf{Hint:} You can use numpy to solve this problem. (5 points)
    \item[(c)] Prove that if a real matrix $\mathbf{A} \in \mathbb{R}^{n \times n}$ has unique eigenvalues, then the eigenvectors $\mathbf{x}_i$ are linearly independent. \textbf{Hint:} Prove by contradiction, start with $n=2$ case. (10 points)
\end{enumerate}

\textbf{Solution:}

\bigskip

%%%%%%%%%%%%%%%%%%%%%%%%%%%%%%%%%%%%%%%%%%%%%%%%%%%%%%%%%%%%%%%%%%%%%%%%
\newpage
\section*{Question 3: Power method (20 points)}

Solve the problem in the \href{https://drive.google.com/file/d/1HQRlmo1zC6ZyeMCt-B278iOXQMfdmMRJ/view?usp=drive_link}{google colab}. Please make a copy of the notebook and solve it there.

\begin{enumerate}
    \item[(a)] Power Method: Write function \texttt{power\_method(A,x)}, which takes as input matrix $A$ and a vector $\mathbf{x}$, and uses the power method to calculate eigenvalues and eigenvectors.

    Get the largest \textbf{(in absolute value)} eigenvalue and the corresponding eigenvector for matrix $A$ using the above function.
    $$
    A=\begin{bmatrix}
        3 & 2 & 1 \\
        1 & 3 & 2 \\
        2 & 4 & 1
    \end{bmatrix}
    $$
    Start with initial eigenvector guesses: $[-1, 0.5, 3]$ and $[2,-6,0.2]$. For each of the vectors, iterate until convergence.
    \begin{enumerate}
        \item[(i)] Plot how the eigenvalue changes with respect to iterations.
        \item[(ii)] Report the number of steps it took to converge, for both the eigenvalue and eigenvector.
        \item[(iii)] Report the final eigenvalue and eigenvector. Match your output with the results generated by the numpy API: \texttt{numpy.linalg.eig}.
    \end{enumerate}
    \emph{Note}: You only need to look at magnitudes of eigenvalues. Use an absolute tolerance of $10^{-6}$ between eigenvalue output of previous and current iteration as stopping criteria. You may also need to normalize the final eigenvector to match with output of numpy API \texttt{numpy.linalg.eig}. (10 points)

    \item[(b)] Inverse Power Method: Write a function \texttt{inverse\_power\_method(A,x)}, which takes as input matrix $\mathbf{A}$ and a vector $\mathbf{x}$, and uses inverse power method to calculate the smallest (in absolute value) eigenvalue and corresponding eigenvector. Solve for the smallest (in absolute value) eigenvalue and corresponding eigenvector for the matrix from (a). \textcolor{blue}{Start with initial eigenvector guesses: $[1, 0.5, 3]$ and $[2,-0.6, 0.2]$. For each of the vectors, iterate until convergence.}
    \begin{enumerate}
        \item[(i)] Plot the computed/estimated eigenvalue with respect to iterations.
        \item[(ii)] Report how many iterations do you need for it to converge to the smallest eigenvalue.
        \item[(iii)] Report the final eigenvalue and eigenvector you get. Match your answer with the results generated by the numpy API \texttt{numpy.linalg.eig}.
    \end{enumerate}
    \emph{Note}: You only need to look at magnitudes of eigenvalues. Use an absolute tolerance of $10^{-6}$ between eigenvalue output of previous and current iteration as stopping criteria. You may also need to normalize the final eigenvector to match with output of numpy API \texttt{numpy.linalg.eig}. (10 points)
\end{enumerate}

\textbf{Solution:}

\bigskip

%%%%%%%%%%%%%%%%%%%%%%%%%%%%%%%%%%%%%%%%%%%%%%%%%%%%%%%%%%%%%%%%%%%%%%%%
\newpage
\section*{Question 4: Face Recognition with Eigenfaces (30 points + 10 Bonus Points)}

Solve the problem in the \href{https://drive.google.com/file/d/1HQRlmo1zC6ZyeMCt-B278iOXQMfdmMRJ/view?usp=drive_link}{google colab}. Please make a copy of the notebook and solve it there.

\textbf{Goal}: Perform face recognition on the \emph{Labeled Faces in the Wild} dataset using PyTorch.

\textbf{Dataset Information}: \href{http://vis-www.cs.umass.edu/lfw/}{Labeled Faces in the Wild} dataset consists of face photographs designed for studying the problem of unconstrained face recognition. The original dataset contains more than 13,000 images of faces collected from the web.

\textbf{Tasks}:
\begin{itemize}
    \item First, perform Principal Component Analysis (PCA) on the image dataset.
    \item Using PCA, extract the top $k$ principal components (eigenvectors). These components, known as \textbf{eigenfaces}, can be used to reconstruct face images.
    \item \textbf{BONUS}: Finally, train a simple PyTorch Neural Network model on the modified image dataset. This trained model will be used for prediction and evaluation on a test set.
\end{itemize}

\textbf{Problem Description}
\begin{enumerate}
    \item[(a)] Preprocessing: Using the \texttt{train\_test\_split} API from \texttt{sklearn}, split the data into train and test dataset in the ratio 4:1. Use \texttt{random\_state=0}.

    For better performance, normalize the features which can have different ranges with huge values.

    Use the \texttt{StandardScaler} class from \texttt{sklearn} and use that to normalize \texttt{X\_train} and \texttt{X\_test}. Validate and show your result by printing the first 5 features of 5 images of \texttt{X\_train} (This result can vary from PC to PC). (10 points)

    \item[(b)] Dimensionality reduction: Use the PCA API from \texttt{sklearn} to extract the top 50 principal components of the image matrix and fit it on the training dataset.

    Visualize some of the top few components as an image (eigenfaces). (10 points)

    \item[(c)] Face reconstruction: Reconstruct an image from its point projected on the principal component basis. Project the first three training images on the eigenvector basis using PCA models trained with varying number of principal components. Using the projected points, reconstruct the faces, and visualize the images.

    Your final output should be a $(3\times 5)$ image matrix, where the rows are the data points, and the columns correspond to original image and reconstructed image for $\text{n\_components} =[10,50,150,500]$. (10 points)

    \item[(d)] Prediction \textbf{(BONUS)}: Train a neural network classifier in \textbf{PyTorch} on the transformed dataset.\\
    Note: For PyTorch reference see \href{https://pytorch.org/docs/stable/index.html}{documentation}. (10 points)
\end{enumerate}

\textbf{Solution:}

\bigskip

\end{document}
